% Responsable: Jerónimo
% Fuentes: docs_proyecto/revision-formulacion-modelo3.md, CLAUDE.md de Modelo3,
% Seleccion_Variables_Restricciones_Modelo3.pdf, desviaciones-paper-m34-6nov.md
% Tesis: el Estado sabe CUÁNTOS privados hay por zona, no QUIÉNES; paga por
% cobertura y el azar decide a quién llega (inclusión/exclusión).
\subsection{Modelo 3: focalización geográfica estocástica}\label{sec:m3}

\subsubsection{Idea}
\pendiente{Gasto por par (indicador, zona) $\to$ cobertura $\cob_{jd}$ $\to$
sorteo $R_{ij}$ $\to$ el hogar sale de la privación si $R_{ij}\le\cob_{jd}$.}

\subsubsection{Formulación}
\pendiente{Objetivo, Con 4*, Con 5*--6*, Con 8--9 y meta de $\Mcero$.}

\subsubsection{Linealización}
\pendiente{Resumen aquí; desarrollo completo en el Apéndice~\ref{app:lin}.}

\subsubsection{Errores de inclusión y exclusión}
\pendiente{Cómo se miden en el modelo y por qué son la ``capacidad real del Estado''.}
